% =====================================================================
%  Método 10 — Plausibilidade do verbo mascarado   (dimensão: semântica)
%  Origem: células 45, 46 e 47 do caderno
% =====================================================================
\subsection{Plausibilidade do verbo mascarado}
\label{met:verbo}

Os métodos anteriores atribuem escore ao que está escrito. Este inverte o procedimento:
suprime-se a palavra e determina-se qual palavra seria esperada naquela posição, confrontando
o texto com a expectativa apenas em seguida.

A supressão de posição arbitrária não seria adequada, e a escolha é sintática: suprime-se a
raiz da árvore de dependências, isto é, o verbo que sustenta a predicação e seleciona os
demais argumentos. O método só se aplica quando $\mathrm{pos}(\raiz(\sent)) = \texttt{VERB}$
e, quando essa condição é falsa --- fragmento ou predicação nominal ---, a resposta adequada é
\texttt{n/a}, e não um escore. Não se trata de limitação implícita, mas de elemento
constitutivo da definição. Cabe observar o arranjo: o analisador sintático determina a posição
a ser consultada, e o modelo estatístico determina o que é plausível naquela posição, sendo
que nenhum dos dois executaria a tarefa isoladamente.

Substituída a raiz pela marca, obtém-se $\sent_{\setminus\raiz}$, e o modelo devolve uma
distribuição sobre todo o vocabulário $\mathcal{V}$, que soma 1. Diferentemente do PLL,
trata-se de probabilidade em sentido estrito, uma vez que a soma percorre as alternativas de
uma única posição, que é exatamente o que o modelo mascarado foi treinado para prever. Dela se
extraem duas leituras, o logaritmo da probabilidade do verbo escrito e a margem até o melhor
preenchimento possível:
%
\begin{equation}
\hat{v} = \arg\max_{v \in \mathcal{V}} P\bigl(v \mid \sent_{\setminus\raiz}\bigr)
\qquad\qquad
m(\sent) = \ln P\bigl(\hat{v} \mid \sent_{\setminus\raiz}\bigr) - \ln P\bigl(v_\raiz \mid \sent_{\setminus\raiz}\bigr)
\end{equation}
%
Por construção $m(\sent) \geq 0$, e $m(\sent) = 0$ exatamente quando o verbo escrito é a
primeira escolha do modelo.

Adota-se a margem, e não o valor absoluto, porque este depende do grau geral de
previsibilidade da posição: em contexto muito restrito toda alternativa recebe probabilidade
elevada e, em contexto vago, toda alternativa recebe probabilidade baixa, inclusive o verbo
adequado. A margem normaliza pelo máximo daquela posição, quantificando a distância ao melhor
preenchimento em lugar da probabilidade absoluta.

O método devolve a margem, acompanhada da lista das formas esperadas, conforme a
Tabela~\ref{tab:verbo}.

\begin{table}[H]
\centering
\dimtable{
\begin{tabular}{@{}lrr@{}}
\toprule
\textbf{Frase} & $\ln P$ (escrito) & \textbf{Margem} \\
\midrule
Há três dias o incêndio \textbf{destruiu} uma área de mata nativa. & $-0{,}72$  & $0{,}00$ \\
Há três dias o incêndio \textbf{atingiu} uma área de mata nativa.  & $-0{,}91$  & $0{,}20$ \\
Há três dias o incêndio \textbf{cantou} uma área de mata nativa.   & $-14{,}45$ & $\mathbf{13{,}73}$ \\
\bottomrule
\end{tabular}}
\caption{A margem do verbo da raiz. As formas esperadas são as mesmas nas três frases:
\texttt{[destruiu, atingiu, atinge, invadiu, em]}.}
\label{tab:verbo}
\end{table}

A separação é ampla, com quase 14 nats entre o verbo adequado e o anômalo e os dois verbos
plausíveis situados junto ao zero. O aspecto metodologicamente relevante é a lista de formas
esperadas ser idêntica nas três frases, o que não poderia ser diferente: ao mascarar a raiz,
suprime-se justamente a única palavra em que as frases diferem, de modo que o contexto
remanescente é o mesmo. O instrumento de medida permanece fixo e apenas o objeto mensurado
varia, e é essa condição que torna a comparação legítima.

A primeira linha exibe o piso da medida. A forma \texttt{destruiu} é a primeira escolha do
modelo, razão pela qual a margem é exatamente nula, o que confere significado ao zero da
escala. Nos casos em que a raiz não é verbo, como no fragmento e na predicação nominal do
corpus, a resposta é \texttt{n/a}, e o reconhecimento da não aplicabilidade integra o emprego
correto do método.
